Introducción:
Se presenta UrbanTracker, un sistema de geolocalización en tiempo real diseñado para optimizar el transporte público urbano mediante el seguimiento continuo de autobuses y la gestión digital de rutas, vehículos y conductores. El proyecto surge ante la falta de información precisa para usuarios y administradores y la necesidad de soluciones asequibles adaptadas a contextos con recursos limitados (p. ej., ciudades intermedias). La validación inicial se realizó en Neiva (Colombia), pero la arquitectura propuesta es transferible a otras ciudades. La motivación se alinea con experiencias previas donde la información en vivo mejora la percepción de puntualidad y la toma de decisiones (“CITY-RUTA”, 2018; “STOPBUS”, 2016).

Métodos:
Se desarrolló una plataforma compuesta por: (i) una app móvil para conductores (GPS del smartphone), (ii) una web pública para usuarios y (iii) un panel web administrativo. El backend sigue una arquitectura hexagonal y modular con separación en domain, application e infrastructure: el modelo de dominio concentra las reglas del negocio; los servicios de aplicación orquestan los casos de uso mediante puertos/contratos; y los adaptadores de infraestructura exponen REST, manejan JPA/PostgreSQL, y conectan con mensajería en tiempo real. Los servicios REST están documentados con OpenAPI/Swagger (EasyRest, 2022), y PostgreSQL gestiona el almacenamiento relacional de rutas, vehículos, usuarios y recorridos.
Para el tiempo real se usó un flujo en el que el móvil del conductor publica posiciones por MQTT hacia un broker; el backend recibe esos mensajes, los valida y guarda, y luego envía actualizaciones por WebSockets a los clientes web (“SisMo”, 2016). La seguridad se maneja con JWT. En la interfaz, la visualización del movimiento en el mapa se implementó con Mapbox en modo 3D.
Para pruebas, se ejecutó un piloto con 1 ruta y 1 vehículo, simulando envíos periódicos de coordenadas cada ~5 s desde el móvil y suscribiendo 5–10 clientes web a las actualizaciones en vivo.

Resultados:
UrbanTracker permite visualizar en el mapa (Mapbox 3D) la ubicación de los autobuses en servicio con una frecuencia de actualización de ~5 segundos, superando un umbral mínimo de 10 s usado como referencia. La latencia medida entre el envío desde el móvil del conductor y la visualización en la web fue < 2 segundos bajo condiciones 4G normales. Se cubren las funcionalidades clave: consulta de rutas, vista de vehículos en tiempo real, inicio/cierre de recorrido por conductores y administración de rutas y flotas desde el panel. En la prueba con 5–10 suscriptores concurrentes, la solución mostró buen comportamiento, en línea con reportes previos del uso de MQTT en rastreo vehicular (“SisMo”, 2016).

Discusión:
La solución cumple lo esperado en el piloto y confirma que es viable brindar información actualizada a pasajeros y herramientas operativas a administradores con infraestructura asequible (smartphones + protocolos ligeros). La arquitectura hexagonal ayuda a separar la captura por MQTT y la difusión por WebSockets, manteniendo responsabilidades claras y facilidades para escalar. Se discuten implicaciones de seguridad y privacidad (uso de JWT y manejo responsable de datos de ubicación) y la compatibilidad con arquitecturas híbridas si crece el histórico de eventos (“Aspectos de ingeniería…”, 2022). La experiencia del usuario fue positiva: el mapa 3D facilita entender dónde está el bus y el panel simplifica tareas de gestión.

Conclusiones:
UrbanTracker evidencia que es factible mejorar la experiencia del transporte público local mediante geolocalización en tiempo real con un stack accesible. Los resultados sugieren impacto potencial en la reducción de tiempos de espera y en la eficiencia operativa. Trabajos futuros incluyen integrar más fuentes de datos en tiempo real, refinar la interfaz y ampliar las evaluaciones de usabilidad a mayor escala, manteniendo la portabilidad del enfoque a otras ciudades.